\chapter{Methodology}
\label{chap:methodology}

\chapterintro{%
This chapter defines the final network-control model and, separately, the exact
contract implemented and replayed on the FPGA. It then specifies candidate fields,
fixed-point scoring, configurations H0--H4, controller thresholds, the H3 rule,
the H4 policy and reward, training partitions, RTL structure, cycle measurement,
physical transport, verification, and statistics. The separation is essential:
the network model describes deployable classical key management, whereas the
completed H4 experiment establishes deterministic controller and hardware behavior
under its registered synthetic inputs.}

\section{Two Related but Distinct Contracts}

The thesis uses two contracts that answer different questions.

\begin{enumerate}
  \item The \emph{network-control model} represents a trusted-relay QKD network
  after link post-processing. It uses classical key-pool occupancy, admitted key
  generation, key consumption, link status, utilization, policy freshness or
  revocation, and externally measured health indicators. It is the intended
  integration model for a KMS/SDN controller.
  \item The \emph{completed H4 experiment} uses the four registered controller
  features: minimum key occupancy, bottleneck route-quality code, offered request
  load, and utilization imbalance. Its second feature retains the registered
  ``fidelity'' field name in source artifacts, but its valid interpretation here is
  a bounded experimental route-quality code.
\end{enumerate}

The second feature of the completed H4 implementation is therefore treated as an
abstract controller input. It is not interpreted as quantum coherence, as fidelity
decay of post-processed keys in classical storage, or as a measured replenishment
rate. Results from that controller support policy-ROM inference, fixed-point
exactness, latency, resource, and trace-specific adaptive behavior. They are not
evidence that the broader network model has already been validated with live QKD
equipment.

\section{System Overview and Boundary}

At a decision epoch, an upstream trusted controller supplies a bounded candidate
set and four state features. The FPGA samples one request, selects a profile,
evaluates the supplied candidates, and returns a registered result. Candidate
generation, graph traversal, authentication, KMS access, security authorization,
and route installation remain outside the kernel.

\begin{figure}[H]
\centering
\resizebox{\textwidth}{!}{%
\begin{tikzpicture}[
  node distance=7mm and 9mm,
  ext/.style={draw=QFlowGray,rounded corners,fill=QFlowLightBlue,
    minimum height=10mm,text width=31mm,align=center,font=\small},
  core/.style={draw=QFlowBlue,rounded corners,fill=white,
    minimum height=10mm,text width=35mm,align=center,font=\small},
  arr/.style={-{Latex[length=2mm]},thick,QFlowNavy}]
\node[ext] (kms) {KMS/SDN monitoring and policy};
\node[ext,right=of kms] (gen) {candidate generation and hard checks};
\node[core,right=of gen] (fpga) {QFlow profile selection and fixed-point evaluation};
\node[ext,right=of fpga] (install) {authorization, installation, fallback};
\draw[arr] (kms)--node[above,font=\scriptsize]{state/metadata}(gen);
\draw[arr] (gen)--node[above,font=\scriptsize]{bounded descriptors}(fpga);
\draw[arr] (fpga)--node[above,font=\scriptsize]{result/status}(install);
\draw[arr,bend left=22] (install) to node[below,font=\scriptsize]{telemetry}(kms);
\end{tikzpicture}}
\caption{System boundary. Only profile selection and candidate evaluation are
included in the measured FPGA kernel.}
\label{fig:system-boundary}
\end{figure}

The kernel-cycle boundary begins on a synchronous edge satisfying
\texttt{start \&\& ready} and ends on the synchronous \texttt{done} edge. UART
serialization, host parsing, operating-system scheduling, USB/JTAG transfer,
programming time, candidate generation, and route installation are excluded.

\section{Final Network and Key-Resource Model}

Let $G=(V,E)$ be a directed trusted-relay QKD network. For link $(i,j)$ at time
$t$, the KMS maintains classical key units $K_{ij}(t)$ with configured capacity
$K_{ij}^{\max}$. A discrete update is
\begin{equation}
K_{ij}[n+1]=\operatorname{clip}_{[0,K_{ij}^{\max}]}
\bigl(K_{ij}[n]+G_{ij}[n]-C_{ij}[n]-X_{ij}[n]\bigr),
\label{eq:method-key-balance}
\end{equation}
where $G_{ij}$ is newly admitted secret material, $C_{ij}$ is consumption, and
$X_{ij}$ is expiry, revocation, or administrative removal. The normalized
occupancy is $o_{ij}=K_{ij}/K_{ij}^{\max}$.

For a request demanding $D$ key units, a link-level hard eligibility indicator is
\begin{equation}
\chi_{ij}(D)=
\mathbf 1\{\mathrm{status}_{ij}=\mathrm{UP}\}
\mathbf 1\{K_{ij}\ge D\}
\mathbf 1\{q_{ij}\le q_{\mathrm{abort}}\}
\mathbf 1\{\mathcal M_{ij}\ \mathrm{valid}\},
\label{eq:method-link-eligibility}
\end{equation}
where $q_{ij}$ is an externally supplied QBER or health measurement and
$\mathcal M_{ij}$ contains authenticated policy metadata, including freshness or
revocation state. Route $\mathcal R_i$ is eligible when
\begin{equation}
v_i=\prod_{(u,v)\in\mathcal R_i}\chi_{uv}(D)=1.
\label{eq:method-route-eligibility}
\end{equation}
These hard conditions are not replaced by a weighted score. Freshness is a
classical KMS policy property; the model contains no stored-key coherence-decay
term.

\section{Executed Request and Candidate Contract}

The physical experiment exposes four candidate slots so that interface width does
not change between builds. Slots 0 and 1 carry the two Ring-6 candidates; slots 2
and 3 are reserved and invalid in the registered replay. Let $N_c$ denote the
number of feasible candidates among the offered slots. $N_c=0$ produces
\texttt{NO\_PATH}; otherwise only feasible candidates enter the objective ranges
and winner reduction.

\begin{landscape}
\begin{table}[p]
\centering
\scriptsize
\setlength{\tabcolsep}{4pt}
\caption{Executed candidate fields used by configurations H0--H4. ``Aggregation''
describes the software/reference source before the fixed-width request is sampled.}
\label{tab:candidate-descriptor}
\begin{tabularx}{\linewidth}{@{}p{2.1cm}p{1.3cm}p{3.2cm}p{2.7cm}p{1.5cm}p{2.0cm}p{2.0cm}X@{}}
\toprule
\textbf{Field} & \textbf{Symbol} & \textbf{Meaning/source} &
\textbf{Route aggregation} & \textbf{Width} & \textbf{Format} &
\textbf{Range} & \textbf{Saturation or exceptional rule}\\
\midrule
Candidate valid & $v_i$ & Upstream eligibility result & Logical AND of link and
request checks & 1 & unsigned Boolean & 0--1 & Invalid candidates never enter
min/max or winner selection\\
Path cost & $c_i$ & Precomputed candidate cost & Sum defined by the software
environment or supplied replay & 32 & unsigned UQ16.16 & $0$--$2^{16}-2^{-16}$ &
Maximum finite code is \texttt{0xFFFFFFFE}; \texttt{0xFFFFFFFF} is infinity\\
Bottleneck route quality & $q_i$ & Registered bounded route-quality input & Minimum
code along the route in the software experiment & 16 & unsigned UNORM16 &
0--65,535, representing 0--1 & Clip to $[0,1]$, multiply by 65,535, round half up\\
Utilization objective & $u_i$ & Candidate utilization/imbalance input & Maximum
candidate-link utilization in the software experiment & 32 & unsigned UQ16.16 &
$0$--$2^{16}-2^{-16}$ & Clip or saturate to the maximum finite code\\
Hop count & $h_i$ & Number of directed edges & Count of edges in route & 3 &
unsigned integer & 0--7 & Values outside the interface width are invalid upstream\\
Candidate slot & $r_i$ & Stable candidate identifier & Assigned by generator & 2 &
unsigned integer & 0--3 & Used as final deterministic tie-break\\
\bottomrule
\end{tabularx}
\end{table}
\end{landscape}

\begin{table}[H]
\centering
\small
\caption{Executed controller-state input fields.}
\label{tab:state-input-fields}
\begin{tabularx}{\textwidth}{@{}p{0.23\textwidth}p{0.16\textwidth}p{0.34\textwidth}p{0.19\textwidth}@{}}
\toprule
\textbf{Field} & \textbf{Width/format} & \textbf{Registered meaning} &
\textbf{Saturation}\\
\midrule
Minimum key occupancy $x_K$ & 16, UNORM16 & Minimum $K/16$ over links appearing
in either supplied Ring-6 candidate & Clip to $[0,1]$\\
Bottleneck route quality $x_Q$ & 16, UNORM16 & Minimum bounded route-quality input
over those links; artifact field name is fidelity & Clip to $[0,1]$\\
Offered request load $x_L$ & 16, UNORM16 & Offered requests in the control window
divided by capacity four & Clip to $[0,1]$\\
Utilization imbalance $x_I$ & 16, UNORM16 & Maximum minus minimum link utilization
over candidate links; 16-step consumption window & Clip to $[0,1]$\\
\bottomrule
\end{tabularx}
\end{table}

All four state ports and all UNORM16 candidate codes use the same conversion:
\begin{equation}
\operatorname{UNORM16}(x)=
\operatorname{clip}_{[0,65535]}
\left\lfloor 65535\operatorname{clip}_{[0,1]}(x)+\tfrac12\right\rfloor.
\label{eq:unorm16}
\end{equation}
Thus one least-significant unit represents $1/65535\approx1.5259\times10^{-5}$
of full scale.

\section{Fixed-Point Objective Formation and Selection}

For each feasible candidate, the three minimization objectives are
\begin{equation}
f_{i,c}=c_i,\qquad f_{i,q}=65535-q_i,\qquad f_{i,u}=u_i.
\label{eq:physical-objectives}
\end{equation}
The route-quality deficit converts a larger-is-better code into a
smaller-is-better objective. For objective $j$, define the ideal point and observed
range over feasible candidates as
\begin{equation}
z_j^\star=\min_{i:v_i=1}f_{i,j},\qquad
F_j^{\max}=\max_{i:v_i=1}f_{i,j}.
\label{eq:ideal-point}
\end{equation}
The ideal point is a componentwise reference and can combine values from different
routes. The normalized penalty is
\begin{equation}
N_{i,j}=\begin{cases}
0,&F_j^{\max}=z_j^\star,\\[2pt]
\operatorname{round}_{\mathrm{half\ up}}
\left(\dfrac{(f_{i,j}-z_j^\star)65535}
{(F_j^{\max}-z_j^\star)+1}\right),&\text{otherwise}.
\end{cases}
\label{eq:method-normalization}
\end{equation}
The $+1$ in the denominator keeps the result finite and inside UNORM16. A single
shared iterative 49-by-33-bit unsigned divider performs the nontrivial range
normalizations.

For profile $\Pi_k$, the 18-bit score is
\begin{equation}
S_i^{(k)}=\max\{\rho_c^{(k)}N_{i,c},
\rho_q^{(k)}N_{i,q},\rho_u^{(k)}N_{i,u}\},
\label{eq:method-score}
\end{equation}
where each $\rho$ is a two-bit unsigned ratio. Lower score is better. The complete
ranking key is
\begin{equation}
\mathcal K_i=(S_i^{(k)},c_i,h_i,r_i),
\label{eq:method-tie-key}
\end{equation}
compared lexicographically in ascending order. Therefore score ties fall back to
raw path cost, then hop count, then candidate slot. If $N_c=0$, score and returned
quality are zero, the selected identifier is 255, and status is
\texttt{NO\_PATH}.

\subsection{Why both alpha and Tchebycheff weights exist}

The software environment can construct path cost from link-level variables using
the registered four-coefficient vector
\begin{equation}
c_i^{(k)}=\sum_{(u,v)\in\mathcal R_i}
\left(\frac{\alpha_1^{(k)}}{K_{uv}}+
\frac{\alpha_2^{(k)}}{Q_{uv}}+
\frac{\alpha_3^{(k)}}{R_{uv}}+
\alpha_4^{(k)}B_{uv}\right),
\label{eq:alpha-path-cost}
\end{equation}
where $Q_{uv}$ is the registered route-quality value, $R_{uv}$ is the synthetic
key-rate variable, and $B_{uv}$ is the synthetic QBER variable. This layer produces
one path-cost objective. The Tchebycheff ratios then trade that objective against
route-quality deficit and utilization imbalance. In the physical comparison,
$c_i$ is supplied precomputed and is common to the replay. The 18-bit profile ROM
retains the $\alpha$ coefficients for semantic/export compatibility, but the
candidate evaluator uses only the lower six ratio bits. Consequently, the physical
adaptive modes reweight three supplied objectives; they do not rerun graph search or
Equation~\eqref{eq:alpha-path-cost} on the FPGA.

\section{Registered Profile Codebook}

\begin{table}[H]
\centering
\scriptsize
\caption{Profile semantics and exact registered coefficients. ``Low-latency''
means path-cost/hop preference; it does not mean shorter FPGA evaluation time.}
\label{tab:profile-codebook}
\begin{tabularx}{\textwidth}{@{}c p{2.4cm} p{3.2cm} p{2.5cm} p{1.8cm} c@{}}
\toprule
\textbf{ID} & \textbf{Registered name} & $\boldsymbol\alpha^{(k)}$ &
$\boldsymbol\lambda_{\rm TCH}^{(k)}$ & \textbf{HW ratio} & \textbf{Payload}\\
\midrule
$\Pi_0$ & Balanced & $(1.0,1.5,0.5,2.0)$ & $(.40,.40,.20)$ & $(2,2,1)$ &
\texttt{13329}\\
$\Pi_1$ & Scarcity protection & $(1.5,1.5,1.0,2.0)$ &
$(.25,.25,.50)$ & $(1,1,2)$ & \texttt{1B516}\\
$\Pi_2$ & Fidelity protection & $(1.0,2.0,0.5,3.0)$ &
$(.25,.50,.25)$ & $(1,2,1)$ & \texttt{14399}\\
$\Pi_3$ & Low-latency & $(0.5,1.0,0.5,1.0)$ &
$(.50,.25,.25)$ & $(2,1,1)$ & \texttt{0A2A5}\\
\bottomrule
\end{tabularx}
\end{table}

The name ``fidelity protection'' is preserved because it is the registered profile
identifier. In this thesis it means stronger emphasis on the bounded route-quality
code in the experimental contract, not protection against coherence decay in a
classical key store. Each $\alpha$ is packed as a three-bit half-unit numerator,
followed by three two-bit Tchebycheff ratios. Multiplying a ratio vector by a common
constant does not alter the selected route, so normalization of $(2,2,1)$ to
$(0.4,0.4,0.2)$ is unnecessary for the integer argmin.

\section{Exact Definitions of Configurations H0--H4}

All five configurations use the same XC6SLX16 target, board shell, clock, candidate
widths, result schema, replay, handshake, and constraint file. They are synthesized,
placed, and routed independently.

\begin{table}[H]
\centering
\small
\caption{Exact selection function of the five physical configurations.}
\label{tab:h0-h4-definitions}
\begin{tabularx}{\textwidth}{@{}c p{0.23\textwidth} X p{0.22\textwidth}@{}}
\toprule
\textbf{Mode} & \textbf{Controller} & \textbf{Winner rule} & \textbf{Purpose}\\
\midrule
H0 & Feasible shortest hop & $\arg\min_{v_i=1}(h_i,r_i)$ & Minimal feasibility and
distance reference\\
H1 & Fixed key-aware cost-only & $\arg\min_{v_i=1}(N_{i,c},c_i,h_i,r_i)$; ratio
$(1,0,0)$ & Cost-only evaluator baseline; no profile-ROM read\\
H2 & Fixed QFlow profile & Equation~\eqref{eq:method-score} with $\Pi_0$ ratio
$(2,2,1)$ & Full fixed-point candidate datapath\\
H3 & Threshold/rule adaptive & Rule in Table~\ref{tab:h3-rules}, followed by the
selected ratio & Deterministic adaptive-controller baseline\\
H4 & Reinforcement-learned adaptive & 256-by-2 policy ROM, three-decision dwell,
then selected ratio & Offline-learned controller under deterministic inference\\
\bottomrule
\end{tabularx}
\end{table}

H0 is deliberately not function-equivalent to H1--H4. H1 evaluates only cost and
is not a full-profile comparator. H2 isolates the complete scoring datapath under a
constant profile. H3 and H4 share the profile-conditioned datapath, making H4 versus
H3 the fairest adaptive-controller comparison.

\begin{pseudocode}{Pseudocode 3.1. Common candidate reduction}
\textbf{Input:} offered candidate slots, selected profile ratio $\boldsymbol\rho$.\\
\textbf{Output:} selected slot, score, quality code, and status.
\begin{enumerate}
  \item Validate request fields and form the feasible set
  $\mathcal F=\{i:v_i=1\}$.
  \item If the request is malformed, return \texttt{INVALID\_REQUEST}.
  \item If $\mathcal F=\varnothing$, return \texttt{NO\_PATH}.
  \item Form cost, quality-deficit, and utilization ranges over $\mathcal F$.
  \item Normalize each objective using Equation~\eqref{eq:method-normalization}.
  \item Compute the applicable H0, H1, or profile score.
  \item Retain the smallest ranking key from Equation~\eqref{eq:method-tie-key}.
  \item Register the winner and assert \texttt{done} for one cycle.
\end{enumerate}
\end{pseudocode}

\section{Controller Quantization and State Encoding}

Each feature is quantized into four levels. For thresholds
$\theta_1<\theta_2<\theta_3$,
\begin{equation}
b(x)=\begin{cases}
0,&x<\theta_1,\\
1,&\theta_1\le x<\theta_2,\\
2,&\theta_2\le x<\theta_3,\\
3,&x\ge\theta_3.
\end{cases}
\label{eq:four-bin}
\end{equation}
Equality always enters the higher bin. Table~\ref{tab:actual-thresholds} gives both
real and encoded thresholds; these are the values compared in the RTL.

\begin{table}[H]
\centering
\scriptsize
\caption{Actual controller thresholds and UNORM16 encodings.}
\label{tab:actual-thresholds}
\begin{tabularx}{\textwidth}{@{}p{2.7cm} c c c p{3.0cm} X@{}}
\toprule
\textbf{Feature} & $\boldsymbol\theta_1$ & $\boldsymbol\theta_2$ &
$\boldsymbol\theta_3$ & \textbf{Encoded thresholds} &
\textbf{Source/meaning}\\
\midrule
Minimum key occupancy & 0.25 & 0.50 & 0.75 & 16,384; 32,768; 49,151 &
Registered capacity-normalized pool thresholds\\
Bottleneck route quality & 0.90 & 0.93 & 0.96 & 58,982; 60,948; 62,914 &
Registered experimental route-quality thresholds\\
Offered request load & 0.25 & 0.50 & 0.75 & 16,384; 32,768; 49,151 &
Requests divided by window capacity four\\
Utilization imbalance & 0.125 & 0.25 & 0.50 & 8,192; 16,384; 32,768 &
Registered max-minus-min utilization thresholds\\
\bottomrule
\end{tabularx}
\end{table}

Let the four bins be $b_K,b_Q,b_L,b_I$. The radix-4 address is
\begin{equation}
z=(((b_K\times4)+b_Q)\times4+b_L)\times4+b_I,
\qquad z\in\{0,\ldots,255\}.
\label{eq:state-address}
\end{equation}
In RTL, $z[7{:}6]=b_K$, $z[5{:}4]=b_Q$, $z[3{:}2]=b_L$, and
$z[1{:}0]=b_I$. This ordering is used by both the H3 rule and H4 policy ROM.

\section{H3 Rule Controller}

H3 applies the first matching row in Table~\ref{tab:h3-rules}. The table is an
executable priority specification, not an informal description.

\begin{table}[H]
\centering
\caption{Exact first-match H3 rule mapping.}
\label{tab:h3-rules}
\begin{tabularx}{\textwidth}{@{}c X c X@{}}
\toprule
\textbf{Priority} & \textbf{Condition} & \textbf{Profile} & \textbf{Numerical
effect in the physical evaluator}\\
\midrule
1 & $b_K=0$ & $\Pi_1$ & Ratio $(1,1,2)$\\
2 & Otherwise, $b_Q\le1$ & $\Pi_2$ & Ratio $(1,2,1)$\\
3 & Otherwise, $b_I\ge2$ & $\Pi_1$ & Ratio $(1,1,2)$\\
4 & Otherwise, $b_L=3\land b_K\ge2\land b_Q\ge2\land b_I\le1$ & $\Pi_3$ &
Ratio $(2,1,1)$\\
5 & Otherwise & $\Pi_0$ & Ratio $(2,2,1)$\\
\bottomrule
\end{tabularx}
\end{table}

For example, a state with both critical occupancy and low route-quality satisfies
rows 1 and 2, but row 1 selects $\Pi_1$. This priority is part of the baseline and
prevents an implementation from reordering conditions for convenience.

\section{H4 Policy-ROM and Dwell Inference}

The H4 policy ROM contains 256 state-major two-bit actions and has one registered
read cycle. The action population is 145 entries selecting $\Pi_0$, 43 selecting
$\Pi_1$, 44 selecting $\Pi_2$, and 24 selecting $\Pi_3$. The profile ROM contains
four 18-bit payloads. The training Q-table contains 1,024 signed 16-bit Q8.7 values
and is retained for analysis only; it is absent from the inference datapath.

The minimum-dwell controller stores the active profile and the number of valid
accepted decisions since the last change. A requested change is admitted only after
the registered dwell of three decisions has been satisfied. A no-path result is a
valid controller decision and advances policy/dwell state; a malformed request is
rejected as invalid and does not become an adaptive decision. Reset clears profile
history and the dwell counter.

\begin{pseudocode}{Pseudocode 3.2. H4 inference}
\textbf{Input:} four UNORM16 state features and current dwell state.\\
\textbf{Output:} deployed profile $a[n]$.
\begin{enumerate}
  \item Quantize $x_K,x_Q,x_L,x_I$ with Table~\ref{tab:actual-thresholds}.
  \item Form address $z[n]$ using Equation~\eqref{eq:state-address}.
  \item Read requested action $a_{\rm req}=\pi(z[n])$ from the policy ROM.
  \item If this is the initial decision, initialize the active action deterministically.
  \item If $a_{\rm req}$ differs from the active action and the dwell requirement is
  not satisfied, retain the active action.
  \item Otherwise admit the requested action and reset the dwell counter on change.
  \item Read the 18-bit profile word and evaluate the candidate set.
\end{enumerate}
\end{pseudocode}

There is no runtime argmax over Q-values, no exploration, no Q update, and no online
learning on the FPGA. ``Reinforcement-learned'' refers to the offline origin of the
ROM mapping.

\section{Immediate Reward and Q-Learning Update}

The original evaluation reward is defined for each decision by success indicator
$y_t$, selected bottleneck route-quality value $b_t\in[0,1]$, post-decision
utilization imbalance $I_t^{\rm post}$, selected hop count $h_t$, action $a_t$, and
previous action $a_{t-1}$. Define
\begin{align}
U_Q(b_t)&=\operatorname{clip}_{[0,1]}
\left(\frac{b_t-0.90}{1-0.90}\right),\\
U_B(I_t^{\rm post})&=1-\operatorname{clip}_{[0,1]}(I_t^{\rm post}),\\
\sigma_t&=\mathbf 1\{a_{t-1}\ \text{exists and }a_t\ne a_{t-1}\}.
\end{align}
For a successful decision,
\begin{equation}
r_t=4+2U_Q(b_t)+U_B(I_t^{\rm post})-0.25h_t-0.25\sigma_t.
\label{eq:success-reward}
\end{equation}
For a blocked decision,
\begin{equation}
r_t=-4-0.25\sigma_t,
\label{eq:block-reward}
\end{equation}
and the quality, balance, and hop components are zero. The declared nominal range
is $[-4.25,6.75]$.

\begin{table}[H]
\centering
\small
\caption{Reward components used by the registered learner and final evaluation.}
\label{tab:reward-components}
\begin{tabularx}{\textwidth}{@{}p{2.1cm}p{1.9cm}p{3.2cm}X@{}}
\toprule
\textbf{Component} & \textbf{Coefficient} & \textbf{Direction and range} &
\textbf{Interpretation/limitation}\\
\midrule
Success & $+4$ & Positive, once per selected route & Dominant service outcome\\
Blocking & $-4$ & Negative, when no route is selected & Does not distinguish the
physical causes of infeasibility\\
Route quality & $+2$ & $0$--$2$ after clipping above 0.90 & Abstract experimental
route-quality utility, not stored-key coherence\\
Balance & $+1$ evaluation; $+4$ training target & Larger when post-decision
imbalance is smaller & Training-target revision changes learning emphasis; reported
evaluation remains on the original $+1$ scale\\
Hop cost & $-0.25$ per edge & More hops reduce reward & A topological proxy, not
measured end-to-end delay\\
Profile switch & $-0.25$ & Applied after the first action when the action changes &
Discourages control chattering; previous action is controller history\\
\bottomrule
\end{tabularx}
\end{table}

The single registered training-target revision changes only the balance coefficient
from 1 to 4 and applies the three-decision dwell during inference; all other reward
terms remain unchanged. Validation and final reported metrics use
Equations~\eqref{eq:success-reward}--\eqref{eq:block-reward}, preserving a common
evaluation scale. This distinction matters because the training target shapes
Q-values, whereas the evaluation reward compares completed controllers.

Tabular Q-learning uses
\begin{equation}
Q(z_t,a_t)\leftarrow Q(z_t,a_t)+0.1
\left[r_t+0.95\max_aQ(z_{t+1},a)-Q(z_t,a_t)\right],
\label{eq:exact-q-update}
\end{equation}
with the terminal bootstrap omitted after the last record. A mean evaluation reward
near 4.67 means that successful decisions dominate and usually earn positive
normalized route-quality and balance utility, partly offset by hop and switch costs.
It is a dimensionless score under this particular reward definition, not money,
energy, secret-key rate, or a universal network utility.

\section{Training Environment and Trace Generation}

\subsection{Topology, candidates, and transition order}

The controlled software environment is a directed Ring-6 with 12 directed edges.
For each source--destination request, the generator supplies the clockwise and
counter-clockwise simple paths. These occupy candidate slots 0 and 1; the four-slot
hardware contract reserves slots 2 and 3. On a successful step, one key unit is
consumed from every link of the selected route. The offered-request count is a
state/load variable in $\{0,1,2,3,4\}$; it is not a claim that a single transition
consumes that many units.

The software step order is:
\begin{enumerate}
  \item observe the current four-feature state;
  \item choose one of four profile actions;
  \item construct the two path objectives and select a route or block;
  \item consume one key unit per selected link on success;
  \item apply the trace's integer key arrivals and registered quality/rate/QBER update;
  \item advance the trace index; and
  \item observe the next state unless the final record was consumed.
\end{enumerate}
No-path is nonterminal. An episode terminates only after record 512. The utilization
feature uses consumption over the preceding 16 steps divided by the current key
count, with denominator at least one.

\subsection{Initialization and phase construction}

Every generated trace begins with 12 independently generated directed-link states.
Initial key count is an integer from 6 through 12 inclusive; initial route-quality
input is 0.9400--0.9900; the synthetic key-rate value is 3--9; and QBER is
0.0150--0.0550. Link capacity is 16 and request-window capacity is four.

Trace phases change every 32 steps and cycle through nominal, scarcity,
route-quality stress (registered as fidelity), high load, low latency, and mixed.
High-load phases set offered load to four; other phases draw an integer from zero
through four using the trace generator's XorShift32 stream. The low-latency phase
changes the source--destination pair from $(0,3)$ to $(0,2)$; other phases use
$(0,3)$. Phase-specific link updates create deterministic differences between the
upper and lower ring directions. These labels describe synthetic trace regimes and
must not be interpreted as observations from deployed quantum links.

\begin{table}[H]
\centering
\small
\caption{Registered trace ranges and their methodological role.}
\label{tab:trace-ranges}
\begin{tabularx}{\textwidth}{@{}p{0.26\textwidth}p{0.25\textwidth}X@{}}
\toprule
\textbf{Quantity} & \textbf{Registered range/rule} & \textbf{Role}\\
\midrule
Initial key count & Integer 6--12; capacity 16 & State occupancy and route
feasibility\\
Initial route-quality input & 0.9400--0.9900 & Abstract bounded controller/path
metric\\
Initial synthetic key rate & 3--9 & Software path-cost and feasibility input\\
Initial QBER & 0.0150--0.0550 & Software path-cost input\\
Offered request load & Integer 0--4; high-load phase fixed at 4 & State feature after
division by four\\
Episode length & 512 decisions & Common finite training/evaluation horizon\\
Phase length & 32 decisions & Repeatable nonstationary regime schedule\\
\bottomrule
\end{tabularx}
\end{table}

In the software environment, a candidate is omitted when any path link has fewer
than one key unit, a route-quality value below 0.90, or nonpositive synthetic key
rate. If neither direction remains, the step blocks but the trace still advances.
The separate physical replay also contains explicitly malformed records to test
\texttt{INVALID\_REQUEST}; those records are protocol tests, not samples from the
training environment.

\subsection{Partitions, seeds, and independence}

\begin{table}[H]
\centering
\small
\caption{Exact trace and trainer seeds. No generated trace appears in more than one
partition.}
\label{tab:training-seeds}
\begin{tabularx}{\textwidth}{@{}p{0.20\textwidth}X@{}}
\toprule
\textbf{Use} & \textbf{Seeds}\\
\midrule
Training traces & 101, 211, 307, 401, 503, 601, 701, 809\\
Validation traces & 1009, 1103, 1201, 1301\\
Test traces & 2003, 2111, 2203, 2309\\
Trainer RNG & 17, 29, 43, 71, 113, 167, 229, 283\\
\bottomrule
\end{tabularx}
\end{table}

Distinct seeds make the generated partitions disjoint and the trainer runs
repeatable, but they do not make all time steps statistically independent: records
within a trace share topology, evolving pools, and phase schedule. Statistical
summaries therefore use seeds or paired physical rows as their declared sampling
units rather than treating every transition as an independent network realization.

\subsection{Training schedule, ranking, and policy freeze}

The Q-table starts at zero. Each trainer candidate runs 200 fixed epochs, eight
episodes per epoch, and 819,200 transitions. Epsilon decreases linearly from 1.0 to
0.05 over the first 80\% of transitions and remains 0.05 thereafter. Evaluation
uses $\epsilon=0$. Each epoch shuffles the eight training traces by XorShift32 and
Fisher--Yates. During exploratory training, equal maximum Q-values are broken
uniformly with XorShift32; validation, testing, and deployment choose the smallest
action identifier on equality.

There is no early stopping or hyperparameter sweep. A candidate must use at least
two profiles, allocate at least 1\% to the second-most-used profile, remain within
an absolute blocking margin of 0.005, and keep switch rate at or below 0.25. The
controlled training revision was permitted once, the test set remained locked until
the revised policy was frozen, and no revision was allowed after test access.
Among seven eligible trainer seeds, seed 229 was selected mechanically by highest
validation reward, followed in order by lower blocking, higher route quality, lower
switching, and smaller seed. The resulting policy and profile images were then used
without online modification.

\section{Detailed RTL Architecture}

\begin{landscape}
\begin{figure}[p]
\centering
\resizebox{0.97\linewidth}{!}{%
\begin{tikzpicture}[
  block/.style={draw=QFlowBlue,rounded corners,fill=white,minimum height=10mm,
    text width=27mm,align=center,font=\scriptsize},
  mem/.style={draw=QFlowTeal,rounded corners,fill=QFlowLightGreen,minimum height=10mm,
    text width=25mm,align=center,font=\scriptsize},
  ctrl/.style={draw=QFlowAmber,rounded corners,fill=QFlowLightAmber,minimum height=10mm,
    text width=27mm,align=center,font=\scriptsize},
  arr/.style={-{Latex[length=2mm]},thick,QFlowNavy},node distance=7mm and 8mm]
\node[block] (in) {input registers\\4 state $\times16$ b\\2 candidates + valid};
\node[block,right=of in] (quant) {four threshold quantizers\\$4\times2$ b};
\node[block,right=of quant] (radix) {radix-4 encoder\\state ID 8 b};
\node[mem,above right=3mm and 9mm of radix] (prom) {H4 policy ROM\\256$\times$2 b};
\node[ctrl,below right=3mm and 9mm of radix] (rule) {H3 first-match\\rule controller};
\node[ctrl,right=18mm of radix] (dwell) {H4 dwell control\\active profile 2 b};
\node[mem,right=of dwell] (prof) {profile ROM\\4$\times$18 b};
\node[block,below=14mm of quant] (store) {candidate storage/stream\\valid, cost, quality, util., hops, slot};
\node[block,right=of store] (feas) {feasibility and range reduction\\min/max per objective};
\node[block,right=of feas] (norm) {range normalizer\\shared 49$\times$33 divider};
\node[block,right=of norm] (score) {ratio products\\Tchebycheff max\\18-b score};
\node[block,right=of score] (win) {winner/tie registers\\score, cost, hops, slot};
\node[ctrl,below=13mm of norm] (fsm) {control FSM\\iteration, bypass, handshakes};
\node[ctrl,right=of fsm] (count) {16-b kernel cycle counter};
\node[block,right=of win] (result) {status/result registers\\policy, state, profile, path, score, quality};
\node[block,below=of result] (fmt) {QF6R formatter + UART\\outside kernel cycles};
\draw[arr] (in)--(quant); \draw[arr] (quant)--(radix);
\draw[arr] (radix)--(prom); \draw[arr] (radix)--(rule);
\draw[arr] (prom)--(dwell); \draw[arr] (rule)--(dwell);
\draw[arr] (dwell)--(prof); \draw[arr] (in)--(store);
\draw[arr] (store)--(feas); \draw[arr] (feas)--(norm);
\draw[arr] (norm)--(score); \draw[arr] (prof)|-(score);
\draw[arr] (score)--(win); \draw[arr] (win)--(result);
\draw[arr] (fsm)-|(store); \draw[arr] (fsm)--(norm);
\draw[arr] (fsm)--(count); \draw[arr] (count)-|(result);
\draw[arr] (result)--(fmt);
\end{tikzpicture}}
\caption{Detailed QFlow RTL architecture. H0 bypasses the profile and iterative
score path; H1 bypasses profile-ROM lookup; H2 fixes $\Pi_0$; H3 uses the rule
controller; H4 uses policy ROM and dwell control. Formatter and UART time are not
included in the kernel-cycle measurement.}
\label{fig:detailed-rtl}
\end{figure}
\end{landscape}

Input registers decouple request sampling from multi-cycle evaluation. State
quantizers and the radix-4 encoder are shared by adaptive modes. The configuration
parameter selects the H0, H1, H2, H3, or H4 controller path at synthesis time. The
candidate engine collects objective extrema, bypasses division when a range is zero,
otherwise iterates the shared divider, forms three weighted penalties, and updates
the winner registers. The FSM sequences all stages and emits exactly one done pulse.
The cycle counter wraps the unchanged policy shell and observes the accepted-edge to
done-edge distance.

\section{Cycle-Level Operation}

A start is accepted on edge $N_{\rm accept}$ only when \texttt{ready} is high. The
wrapper zeroes its counter on that edge. Initialization clears range and winner
registers; candidate collection records feasibility and extrema; the evaluator then
normalizes nonconstant objective ranges, computes the score, compares the winner,
and registers the result. On the one-cycle done pulse,
\begin{equation}
L=N_{\rm done}-N_{\rm accept},\qquad t=L(10~\mathrm{ns})
\label{eq:latency}
\end{equation}
at the physical 100~MHz clock.

\begin{table}[H]
\centering
\small
\caption{Observed normal-result cycle structure. The min/max/median values are
measured from the common replay, not estimates from a nominal pipeline formula.}
\label{tab:cycle-operation}
\begin{tabularx}{\textwidth}{@{}c rrrr X@{}}
\toprule
\textbf{Mode} & \textbf{Min} & \textbf{Mean} & \textbf{Median} & \textbf{Max} &
\textbf{Engineering explanation}\\
\midrule
H0 & 2 & 2.000 & 2 & 2 & Direct feasible-hop comparison followed by
registered completion; iterative evaluator bypassed\\
H1 & 7 & 201.763 & 309 & 309 & Cost-only range normalization; profile-ROM
read bypassed\\
H2 & 8 & 202.763 & 310 & 310 & H1-style evaluation plus one profile-ROM
setup cycle for fixed $\Pi_0$\\
H3 & 8 & 202.763 & 310 & 310 & Rule selection completes within the same
setup envelope as H2\\
H4 & 10 & 204.763 & 312 & 312 & Policy-ROM and dwell/controller pipeline add
two cycles relative to H3\\
\bottomrule
\end{tabularx}
\end{table}

The shared iterative divider dominates H1--H4. Equal objective ranges and
single-feasible-candidate cases bypass some iterations, which explains the small
minimum but much larger median and the noninteger mean over 76 normal rows. H0 is
roughly two cycles because it never enters this normalization path. H1 is one cycle
shorter than H2/H3 because its cost-only controller does not read the profile ROM.
H4 adds exactly two setup cycles to the otherwise comparable H3 datapath. Thus H4
versus H3 is the appropriate controller-overhead comparison; H4 versus H0 mixes a
controller change with the addition of the complete multi-objective evaluator.

\section{Physical Result Protocol}

The synchronous shell returns the nine fields in Table~\ref{tab:nine-result-fields}.
Every one was compared with the software reference for every physical record.

\begin{table}[H]
\centering
\scriptsize
\caption{Nine exact fields in the physical replay result.}
\label{tab:nine-result-fields}
\begin{tabularx}{\textwidth}{@{}p{2.7cm}p{1.4cm}p{2.4cm}X@{}}
\toprule
\textbf{Field} & \textbf{Decimal width} & \textbf{Range} & \textbf{Meaning}\\
\midrule
\texttt{test\_id} & 4 & 0--9999 & Stable replay record identifier\\
\texttt{policy\_id} & 1 & 0--4 & Physical configuration H0--H4\\
\texttt{state\_id} & 3 & 0--255 & Encoded controller state; zero for invalid request\\
\texttt{profile\_id} & 1 & 0--3 & Applied profile; zero for invalid request\\
\texttt{selected\_path} & 3 & 0--3 or 255 & Candidate slot, or 255 on failure\\
\texttt{score} & 6 & 0--262143 & Unsigned 18-bit winner score; zero on failure\\
\makecell[l]{\texttt{bottleneck}\\\texttt{\_fidelity}} & 5 & 0--65535 & Registered name for returned route-quality
code; zero on failure\\
\texttt{cycles} & 3 & 0--999 & $N_{\rm done}-N_{\rm accept}$\\
\texttt{status} & 1 & 0--2 & Completion status\\
\bottomrule
\end{tabularx}
\end{table}

Status 0 is \texttt{OK}, status 1 is \texttt{NO\_PATH}, and status 2 is
\texttt{INVALID\_REQUEST}. For status 1 or 2 the selected path is 255 and score and
returned quality are zero. An accepted invalid request completes in one kernel
cycle. A start presented while \texttt{busy} is high is not accepted, does not reset
the counter, produces no separate response, and cannot disturb the active request.

The transport record is exactly
\begin{center}
\texttt{QF6R,1,TTTT,P,SSS,R,PPP,CCCCCC,FFFFF,LLL,E\textbackslash r\textbackslash n}
\end{center}
and is 44 ASCII bytes including CRLF. Fields are unsigned decimal with intentional
leading zeroes and are emitted left to right in the order shown. The production
link is 115,200 baud, 8 data bits, no parity, one stop bit, idle high, and no flow
control. Each UART frame sends a low start bit, data bits least-significant-bit
first, and a high stop bit. The rounded divider is 868 clocks per serial bit. The
formatter uses 18 iterative shift-add-3 steps per numeric field; formatter and UART
latency are outside Equation~\eqref{eq:latency}.

\section{Verification Methodology}

Verification proceeds from pure functions to physical records. A higher level is
entered only after the lower-level numeric and interface contracts pass.

\begin{figure}[H]
\centering
\begin{tikzpicture}[
  level/.style={draw=QFlowBlue,rounded corners,minimum width=10.8cm,
    minimum height=9mm,align=center,font=\small},
  arr/.style={-{Latex[length=2mm]},thick,QFlowNavy},node distance=5mm]
\node[level,fill=QFlowLightBlue] (l1) {1. Python unit tests: quantization, profile
packing, reward, route selection};
\node[level,fill=QFlowLightBlue,above=of l1] (l2) {2. Exhaustive controller tests:
256 states, four actions, threshold edges, ROM exactness};
\node[level,fill=QFlowLightGreen,above=of l2] (l3) {3. RTL unit tests: encoder,
policy/profile ROMs, H3 rule, dwell, evaluator, protocol};
\node[level,fill=QFlowLightGreen,above=of l3] (l4) {4. Integrated RTL replay:
fixed vectors, invalid/no-path, reset, busy, cycle counts};
\node[level,fill=QFlowLightAmber,above=of l4] (l5) {5. Independent synthesis,
map, place-and-route, timing and bitstream generation for H0--H4};
\node[level,fill=QFlowLightAmber,above=of l5] (l6) {6. Same-board UART replay:
raw capture, parse, nine-field exact comparison, cross-mode statistics};
\draw[arr] (l1)--(l2); \draw[arr] (l2)--(l3); \draw[arr] (l3)--(l4);
\draw[arr] (l4)--(l5); \draw[arr] (l5)--(l6);
\end{tikzpicture}
\caption{Verification ladder from executable arithmetic to same-board physical
capture. Each level tests a distinct failure class.}
\label{fig:verification-ladder}
\end{figure}

The software suite includes 74 regression tests. Exhaustive tests cover all 256
state addresses, all four actions, all policy and profile ROM words, and below/at/
above cases for every quantizer threshold. RTL tests also cover both candidates,
ties, equal objective ranges, one feasible candidate, no path, invalid request,
reset, back-to-back acceptance, and a start while busy. Integrated tracing retained
2,125 cycle vectors and 530 valid decisions. Physical, post-route, RTL, and software
evidence are reported separately.

\section{FPGA Implementation and Physical Replay}

The target is a Digilent Nexys 3 Rev. B containing a Xilinx
XC6SLX16-2-CSG324 Spartan-6 FPGA. Xilinx ISE 14.7 independently builds H0--H4. The
physical clock is 100~MHz; post-route $f_{\max}$ is implementation capacity and is
not substituted for the board clock when converting cycles to time.

The common corpus contains 20 contract rows and 64 held-out dynamic rows, for 84
records per configuration. For each build, the host programs the volatile
bitstream, transmits one registered request at a time, stores the raw UART capture,
parses the nine output fields, and compares them with the fixed-point reference.
The experiment therefore contains 420 physical records and 3,780 field
comparisons. The held-out partition was materialized with seed 26072026 before the
cross-mode result tables were generated.

\begin{pseudocode}{Pseudocode 3.3. Same-board replay and exact comparison}
\textbf{Input:} independently built configuration $h$, common replay $\mathcal D$,
expected records $Y_h$.\\
\textbf{Output:} raw capture, parsed table, mismatch report.
\begin{enumerate}
  \item Program the bitstream for $h$ and verify the expected configuration identity.
  \item Present each replay request only when the shell is ready.
  \item Retain the raw 44-byte response before parsing.
  \item Parse all nine fields and reject missing or duplicate test identifiers.
  \item Compare every integer field with $Y_h$; do not use numeric tolerances.
  \item Retain cycle, status, profile, and selected-route summaries only after exactness passes.
\end{enumerate}
\end{pseudocode}

A board photograph was not archived in the registered repository and is therefore
not introduced as evidence. Figure~\ref{fig:system-boundary} and the transport
specification describe the setup that can be reproduced from retained files.

\section{Statistical Analysis}

\subsection{Training-level summaries}

Controller evaluation uses four validation or test trace seeds as the paired unit.
The registered training-level analysis uses 10,000 paired-seed percentile bootstrap
replicates with RNG seed 32452843, exact paired permutation tests, and Holm
adjustment for the planned comparisons. Final descriptive results are reported on
the original reward scale. With four pairs, the smallest attainable two-sided exact
permutation $p$-value is coarse; the reported raw value 0.125 and Holm-adjusted value
0.25 are therefore interpreted descriptively rather than as proof of superiority.

\subsection{Physical route-quality summaries}

The physical paired sampling unit is one jointly successful held-out replay row.
For comparison $A-B$, difference $d_i=q_{i,A}-q_{i,B}$ is in UNORM16 code units, so
positive means higher route quality in $A$. The paired standardized effect is
\begin{equation}
d_z=\frac{\bar d}{s_d},
\label{eq:paired-dz}
\end{equation}
where $s_d$ is the sample standard deviation of nonaggregated paired differences.
The exact two-sided sign test discards zero differences and applies a binomial test
to the remaining positive and negative signs. The two physical comparisons are
treated as exploratory and their $p$-values are not multiplicity-adjusted.

The retained physical summary provides percentile bootstrap interval endpoints and
identifies jointly successful rows as the resampling unit, but it does not preserve
the bootstrap replication count or RNG seed. Those quantities cannot be recovered
without inventing provenance. Accordingly, the thesis reports the registered
intervals, does not claim fully reproducible resampling internals for this one
analysis, and lists the omission as a statistical reproducibility limitation.
Metric direction, units, pair counts, exact sign-test behavior, and effect-size
formula remain fully specified.

\section{Supporting VLSI Methodology}

The supplementary VLSI study uses isolated arithmetic and comparison kernels. Yosys
generic synthesis compares coefficient-specialized and general variants. Selected
kernels are separately processed through a SKY130 high-density OpenROAD flow with
synthesis, floorplanning, placement, clock-tree synthesis, routing, parasitic
extraction, and report generation. Generic cells are not FPGA LUTs, and isolated
kernel areas are not summed into a QFlow ASIC area.

The repository identifier FDPE-V3 is described in the thesis as a
``LUT/interpolation arithmetic kernel (repository identifier FDPE-V3).'' It is a
supplementary numerical experiment, not a fidelity-decay function in the current
QFlow candidate evaluator. A representative 2:1 transmission-gate multiplexer is
also simulated in ngspice with simplified Level-1 models at 1.8~V and 5--50~fF
loads. The calculated $\tfrac12CV^2$ quantity applies only to the stated capacitive
load and is not accelerator energy per decision.

\section{Methodological Limitations}

The controlled learning environment uses one Ring-6 topology, synthetic phases, two
valid candidates per request, and a 512-step horizon. The physical experiment uses
one board/device family and 84 records per configuration. The second controller
feature is an abstract bounded route-quality input, not a live-link measurement or
stored-key quantum state. No matched CPU implementation, activity-calibrated FPGA
power measurement, live QKD devices, authenticated KMS/SDN deployment, or installed
network route timing was completed. Physical bootstrap replication metadata is
incomplete. The VLSI study covers isolated kernels and is not full-chip signoff.

\chapterconclusion{%
The methodology now specifies the exact executed semantics: four registered
controller features, numerical profile codebook, explicit H0--H4 selection laws,
actual thresholds, complete reward, trace partitions, deterministic ROM inference,
fixed-point arithmetic, cycle boundary, nine-field UART record, and statistical
units. It also keeps those hardware results separate from the deployable classical
key-resource model, preventing an abstract experimental route-quality code from
being presented as stored-key coherence or replenishment evidence.}
